% titel: Blatt 1, A4: Polarkoordinaten
% skript: Kap. Kinematik, ebene Polarkoordinaten

% Einheitsvektoren, Orts- und Geschwindigkeitsvektor in Polarkoordinaten.

Bei ebenen Bahnen bietet ein Koordinatensystem an den Ort die
Polarkoordinaten $(r,\varphi)$: Grundvektoren heissen $\hat e_r$
(radial) und $\hat e_\varphi$ (tangential zur Kreisbahn).

\begin{werkzeugkasten}[Werkzeuge]
\textbf{Umrechnung} kartesisch $\leftrightarrow$ polar:
\begin{equation}
  x=r\cos\varphi,\quad y=r\sin\varphi,\qquad
  r=\sqrt{x^2+y^2}.
\end{equation}

\textbf{Einheitsvektoren} (normierte partielle Ableitungen):
\begin{equation}
  \hat e_r=\frac{\partial\vec r/\partial r}{|\cdot|},\qquad
  \hat e_\varphi=\frac{\partial\vec r/\partial\varphi}{|\cdot|}\;.
\end{equation}

\textbf{Ortsvektor} in Polarkoordinaten:
\begin{equation}
  \vec r(\vec r,\varphi)=r\,\hat e_r\;.
\end{equation}

\textbf{Zeitableitung von $\hat e_r$ und $\hat e_\varphi$:} Weil sich die
Basis unter der Drehung $\varphi$ ändert, gilt
\begin{equation}
  \dot{\hat e}_r=\dot\varphi\,\hat e_\varphi,\qquad
  \dot{\hat e}_\varphi=-\dot\varphi\,\hat e_r\;.
\end{equation}

\textbf{Geschwindigkeit} (Produktregel $\vec r=r\hat e_r$):
\begin{equation}
  \vec v=\dot r\,\hat e_r+r\,\dot\varphi\,\hat e_\varphi\;.
\end{equation}
\end{werkzeugkasten}

\begin{werkzeugkasten}[Vorgehen]
\begin{itemize}
  \item Einheitsvektoren: partielle Ableitung bilden, dann normieren.
  \item Die Zeitableitung der Basis nicht vergessen: die Einheitsvektoren
        sind orts- (und damit zeit-)abhängig.
  \item Bei der Kreisbewegung ist $\dot r=0$ und $\varphi=\omega t$; der
        Radialteil verschwindet.
\end{itemize}
\end{werkzeugkasten}

\begin{achtung}[Typische Fallen]
\begin{itemize}
  \item $\hat e_r$ und $\hat e_\varphi$ sind \emph{nicht} konstant, anders
        als $\hat e_x,\hat e_y$.
  \item Beim Ableiten von $r\hat e_r$ die Produktregel verwenden; sonst
        fehlt der tangentiale Term.
  \item Für die Kreisbahn mit $r=R$ die Länge $|\vec v|$ aus beiden
        Komponenten ausrechnen, nicht nur den Betrag von $\hat e_\varphi$
        (der ist $1$).
\end{itemize}
\end{achtung}
